% chapters/02/evaluasi.tex
\bagian{Asesmen dan Instrumen Penilaian}

\begin{asesmenbox}
\textbf{Kedudukan asesmen bab ini.} Bab~II juga tidak memiliki Sub-CPMK
tersendiri karena berperan sebagai landasan teori bagi seluruh bab berikutnya.
Asesmennya bersifat \textbf{formatif} dan tidak diberi bobot pada nilai akhir.
Instrumen di bawah ini dirancang untuk menemukan salah pengertian
(\istilah{miskonsepsi}) yang apabila dibiarkan akan menyulitkan Anda pada
praktikum Bab~III dan seterusnya. Jawablah tanpa membuka bab terlebih dahulu,
lalu periksa jawaban Anda dengan kunci yang tersedia.
\end{asesmenbox}

\subsection*{Instrumen: Pemeriksaan Salah Pengertian}

Tentukan benar atau salah setiap pernyataan berikut, lalu tuliskan alasan
singkat atas pilihan Anda.

\begin{enumerate}
  \item Memperbanyak jumlah bingkai per detik selalu membuat permainan menjadi
        lebih baik.
  \item Kecepatan objek sebaiknya dinyatakan dalam satuan per bingkai agar
        perhitungannya lebih sederhana.
  \item Pada koordinat layar, menambahkan nilai pada koordinat tegak berarti
        objek bergerak ke bawah.
  \item Kaidah Euler menghasilkan lintasan yang sama persis dengan lintasan
        yang sebenarnya.
  \item Memperkecil panjang langkah memperkecil galat kaidah Euler.
  \item Pada mesin keadaan, NPC diperbolehkan berada pada dua keadaan
        sekaligus.
  \item Dua objek yang persegi pembatasnya beririsan pasti gambar keduanya
        bersinggungan.
  \item Pada mesin permainan, alur program utama tetap ditulis sendiri oleh
        pembuat permainan.
  \item Menyimpan satu informasi keadaan yang sama di dua tempat berbeda
        memudahkan pemeliharaan program.
  \item Pemahaman tentang cara kerja siklus permainan tidak lagi diperlukan
        setelah beralih memakai mesin permainan.
\end{enumerate}

\begin{table}[htbp]
\centering
\caption{Kunci dan penjelasan pemeriksaan salah pengertian}
\label{tab:kunci-miskonsepsi}
{\setlength{\tabcolsep}{4pt}\footnotesize%
\begin{tabularx}{\textwidth}{@{}C{0.8cm} C{1.2cm} Y@{}}
\toprule
\headrow \textbf{No} & \textbf{Kunci} & \textbf{Penjelasan} \\
\midrule
1 & Salah &
Laju bingkai yang lebih tinggi menuntut anggaran waktu yang lebih ketat dan
memakan daya lebih besar. Yang menentukan mutu bukan angka bingkai, melainkan
konsistensi perilaku permainan. \\
\addlinespace
2 & Salah &
Kecepatan dalam satuan per bingkai membuat gerak bergantung pada kecepatan
komputer, sehingga permainan menjadi tidak adil dan tidak dapat diramalkan. \\
\addlinespace
3 & Benar &
Koordinat layar ber titik asal di kiri-atas dengan sumbu tegak menuju ke bawah,
sehingga penambahan nilai berarti perpindahan ke bawah. \\
\addlinespace
4 & Salah &
Kaidah Euler adalah hampiran. Selisihnya menyusut ketika panjang langkah
diperkecil, tetapi tidak pernah menjadi nol. \\
\addlinespace
5 & Benar &
Galat kaidah Euler berbanding lurus dengan panjang langkah. \\
\addlinespace
6 & Salah &
Batasan ``tepat satu keadaan pada satu waktu'' justru menjadi kekuatan
pendekatan ini karena menghapus kebutuhan memikirkan kombinasi keadaan. \\
\addlinespace
7 & Salah &
Persegi pembatas hanyalah pendekatan. Dua objek berlubang dapat terdeteksi
bersentuhan walaupun gambarnya tidak bersinggungan. \\
\addlinespace
8 & Salah &
Mesin memiliki alur programnya sendiri dan memanggil kode Anda pada tahap yang
telah ditentukan. \\
\addlinespace
9 & Salah &
Dua salinan informasi yang sama cenderung saling bertentangan pada suatu saat.
Keadaan sebaiknya dinyatakan terpusat dan sejelas mungkin. \\
\addlinespace
10 & Salah &
Mesin menyembunyikan sebab akibat, sehingga pemahaman tentang mekanisme dasar
justru menjadi bekal untuk mengenali kapan pengaturan bawaan perlu diubah. \\
\bottomrule
\end{tabularx}}
\end{table}

\begin{catatabox}
Apabila Anda menjawab salah pada butir nomor 2, 4, atau 5, bacalah kembali
Bagian~\ref{sec:model-waktu} dan Bagian~\ref{sec:simulasi-diskret} sebelum
memulai Bab~III. Ketiga butir itu menyangkut hal yang paling sering menjadi
sumber kesalahan pada praktikum Blok~I.
\end{catatabox}
